% --- Archon physics formalization source begin ---
% archon:physics
% archon:covers PhyXMiniProblems/problem_phyx_mini_0571.lean
% archon:source-report reports/phyx_mini/problem_phyx_mini_0571.source.json

\chapter{Physics problem phyx\_mini\_0571}
\label{ch:PhyXMiniProblems_problem_phyx_mini_0571}

\paragraph{Problem source.}
\#\# Physical scenario

Consider the model of the H\$\_2\$O molecule shown in the diagram.

\#\# Question

What is the wavelength of a photon required to excite a transition from \$l=0\$ to \$l=1\$?

\#\# Answer choices

- A: A: 123\textbackslash{}mu\textbackslash{}m

- B: B: 452\textbackslash{}mu\textbackslash{}m

- C: C: 344\textbackslash{}mu\textbackslash{}m

- D: D: 425\textbackslash{}mu\textbackslash{}m

\#\# Figure caption (auxiliary; use the image as primary evidence)

The image depicts a molecular structure of a water molecule. It consists of three atoms: two hydrogen (H) atoms and one oxygen (O) atom. 

- The oxygen atom is larger and centrally located, shaded in a dark gray color.
- The two hydrogen atoms are smaller and positioned on either side of the oxygen atom, shaded in a light blue color.
- The hydrogen atoms are connected to the oxygen atom by lines representing chemical bonds.
- The angle between the two hydrogen atoms, with oxygen as the vertex, is labeled as 105°.
- The length of each bond between the hydrogen and oxygen atoms is labeled as 0.0958 nm.
- A dashed line is shown to indicate the reference line for the angle measurement.
  
This is a common representation of the molecular geometry of water, showing its bent shape.

\#\# Dataset metadata

Quantum Phenomena; Physical Model Grounding Reasoning, Numerical Reasoning

\paragraph{Recorded answer/context.}
C: C: 344\textbackslash{}mu\textbackslash{}m

\paragraph{Figure/image path.}
/root/proposal\_for\_physic/science-mango/phyx\_mini\_run/phyx\_data/test\_image/571.png

\paragraph{Formalization target.}
create a compiling Lean file with sorry bodies at `PhyXMiniProblems/problem\_phyx\_mini\_0571.lean`. The Lean declarations must preserve the physical quantities, dimensions or dimensional roles, figure labels, governing-law hypotheses, and final relation expressed by this problem.
Use Mathlib/Physlib names found through LeanExplore where available. If a physics API is missing, introduce faithful local abstractions rather than scalar placeholder aliases.

\begin{theorem}[Physics formalization target]
\label{thm:physics:phyx_mini_0571:target}
\lean{PhyXMiniProblems.ProblemPhyXMini0571.requiredPhotonWavelength_is_choiceC}
\uses{def:physics:phyx-mini-0571:phyxminiproblems-problemphyxmini0571-watermoleculerotorsetup, def:physics:phyx-mini-0571:phyxminiproblems-problemphyxmini0571-matcheswaterrotationalexcitationscenario, def:physics:phyx-mini-0571:phyxminiproblems-problemphyxmini0571-matchesprimarywatermoleculefigure, def:physics:phyx-mini-0571:phyxminiproblems-problemphyxmini0571-usestextbookreferencedata, def:physics:phyx-mini-0571:phyxminiproblems-problemphyxmini0571-hasphysicalwaterrotorparameters, def:physics:phyx-mini-0571:phyxminiproblems-problemphyxmini0571-satisfiesdepictedbisectorgeometry, def:physics:phyx-mini-0571:phyxminiproblems-problemphyxmini0571-satisfiesrigidrotorphotonlaws, def:physics:phyx-mini-0571:phyxminiproblems-problemphyxmini0571-agreeswithdisplayedwavelength, def:physics:phyx-mini-0571:phyxminiproblems-problemphyxmini0571-isuniqueclosestwavelengthchoice}
The assigned autoformalize agent should translate this physics problem into Lean declarations in the covered file, with theorem and lemma proof bodies written as `by sorry`.
\end{theorem}
\begin{proof}
This is an autoformalization task, not a proof task. Produce statements that can later be proved by the physics prover without weakening the physical model.
\end{proof}

% --- Archon declaration topology begin ---
\section{Declaration topology}
\begin{definition}[Length Quantity]
  \label{def:physics:phyx-mini-0571:phyxminiproblems-problemphyxmini0571-lengthquantity}
  \lean{PhyXMiniProblems.ProblemPhyXMini0571.LengthQuantity}
  A nonnegative, unit-independent physical length
\end{definition}
\begin{proof}
  This is the declared model definition of Length Quantity.
\end{proof}
\begin{definition}[Mass Quantity]
  \label{def:physics:phyx-mini-0571:phyxminiproblems-problemphyxmini0571-massquantity}
  \lean{PhyXMiniProblems.ProblemPhyXMini0571.MassQuantity}
  A nonnegative, unit-independent physical mass
\end{definition}
\begin{proof}
  This is the declared model definition of Mass Quantity.
\end{proof}
\begin{definition}[moment Of Inertia Dimension]
  \label{def:physics:phyx-mini-0571:phyxminiproblems-problemphyxmini0571-momentofinertiadimension}
  \lean{PhyXMiniProblems.ProblemPhyXMini0571.momentOfInertiaDimension}
  The physical dimension of a moment of inertia, mass * length\textasciicircum{}2
\end{definition}
\begin{proof}
  This is the declared model definition of moment Of Inertia Dimension.
\end{proof}
\begin{definition}[Moment Of Inertia Quantity]
  \label{def:physics:phyx-mini-0571:phyxminiproblems-problemphyxmini0571-momentofinertiaquantity}
  \lean{PhyXMiniProblems.ProblemPhyXMini0571.MomentOfInertiaQuantity}
  \uses{def:physics:phyx-mini-0571:phyxminiproblems-problemphyxmini0571-momentofinertiadimension}
  A nonnegative, unit-independent moment of inertia
\end{definition}
\begin{proof}
  This is the declared model definition of Moment Of Inertia Quantity.
\end{proof}
\begin{definition}[action Dimension]
  \label{def:physics:phyx-mini-0571:phyxminiproblems-problemphyxmini0571-actiondimension}
  \lean{PhyXMiniProblems.ProblemPhyXMini0571.actionDimension}
  The physical dimension of action, mass * length\textasciicircum{}2 / time
\end{definition}
\begin{proof}
  This is the declared model definition of action Dimension.
\end{proof}
\begin{definition}[Action Quantity]
  \label{def:physics:phyx-mini-0571:phyxminiproblems-problemphyxmini0571-actionquantity}
  \lean{PhyXMiniProblems.ProblemPhyXMini0571.ActionQuantity}
  \uses{def:physics:phyx-mini-0571:phyxminiproblems-problemphyxmini0571-actiondimension}
  A nonnegative, unit-independent physical action
\end{definition}
\begin{proof}
  This is the declared model definition of Action Quantity.
\end{proof}
\begin{definition}[Speed Quantity]
  \label{def:physics:phyx-mini-0571:phyxminiproblems-problemphyxmini0571-speedquantity}
  \lean{PhyXMiniProblems.ProblemPhyXMini0571.SpeedQuantity}
  A unit-independent speed with a real carrier. This matches the exact type of Physlib's dimensionful DimSpeed.speedOfLight constant
\end{definition}
\begin{proof}
  This is the declared model definition of Speed Quantity.
\end{proof}
\begin{definition}[length Readout]
  \label{def:physics:phyx-mini-0571:phyxminiproblems-problemphyxmini0571-lengthreadout}
  \lean{PhyXMiniProblems.ProblemPhyXMini0571.lengthReadout}
  \uses{def:physics:phyx-mini-0571:phyxminiproblems-problemphyxmini0571-lengthquantity}
  Read a physical length as a real number in a selected length unit
\end{definition}
\begin{proof}
  This is the declared model definition of length Readout.
\end{proof}
\begin{definition}[length In Meters]
  \label{def:physics:phyx-mini-0571:phyxminiproblems-problemphyxmini0571-lengthinmeters}
  \lean{PhyXMiniProblems.ProblemPhyXMini0571.lengthInMeters}
  \uses{def:physics:phyx-mini-0571:phyxminiproblems-problemphyxmini0571-lengthquantity, def:physics:phyx-mini-0571:phyxminiproblems-problemphyxmini0571-lengthreadout}
  Metre readout of a physical length
\end{definition}
\begin{proof}
  This is the declared model definition of length In Meters.
\end{proof}
\begin{definition}[length In Nanometers]
  \label{def:physics:phyx-mini-0571:phyxminiproblems-problemphyxmini0571-lengthinnanometers}
  \lean{PhyXMiniProblems.ProblemPhyXMini0571.lengthInNanometers}
  \uses{def:physics:phyx-mini-0571:phyxminiproblems-problemphyxmini0571-lengthquantity, def:physics:phyx-mini-0571:phyxminiproblems-problemphyxmini0571-lengthreadout}
  Nanometre readout of a physical length
\end{definition}
\begin{proof}
  This is the declared model definition of length In Nanometers.
\end{proof}
\begin{definition}[length In Micrometers]
  \label{def:physics:phyx-mini-0571:phyxminiproblems-problemphyxmini0571-lengthinmicrometers}
  \lean{PhyXMiniProblems.ProblemPhyXMini0571.lengthInMicrometers}
  \uses{def:physics:phyx-mini-0571:phyxminiproblems-problemphyxmini0571-lengthquantity, def:physics:phyx-mini-0571:phyxminiproblems-problemphyxmini0571-lengthreadout}
  Micrometre readout of a physical length
\end{definition}
\begin{proof}
  This is the declared model definition of length In Micrometers.
\end{proof}
\begin{definition}[mass In Kilograms]
  \label{def:physics:phyx-mini-0571:phyxminiproblems-problemphyxmini0571-massinkilograms}
  \lean{PhyXMiniProblems.ProblemPhyXMini0571.massInKilograms}
  \uses{def:physics:phyx-mini-0571:phyxminiproblems-problemphyxmini0571-massquantity}
  Kilogram readout of a physical mass
\end{definition}
\begin{proof}
  This is the declared model definition of mass In Kilograms.
\end{proof}
\begin{definition}[moment Of Inertia In Kilogram Meters Squared]
  \label{def:physics:phyx-mini-0571:phyxminiproblems-problemphyxmini0571-momentofinertiainkilogrammeterssquared}
  \lean{PhyXMiniProblems.ProblemPhyXMini0571.momentOfInertiaInKilogramMetersSquared}
  \uses{def:physics:phyx-mini-0571:phyxminiproblems-problemphyxmini0571-momentofinertiaquantity}
  Coherent-SI readout of a moment of inertia, in kg m²
\end{definition}
\begin{proof}
  This is the declared model definition of moment Of Inertia In Kilogram Meters Squared.
\end{proof}
\begin{definition}[action In Joule Seconds]
  \label{def:physics:phyx-mini-0571:phyxminiproblems-problemphyxmini0571-actioninjouleseconds}
  \lean{PhyXMiniProblems.ProblemPhyXMini0571.actionInJouleSeconds}
  \uses{def:physics:phyx-mini-0571:phyxminiproblems-problemphyxmini0571-actionquantity}
  Coherent-SI readout of an action, in joule-seconds
\end{definition}
\begin{proof}
  This is the declared model definition of action In Joule Seconds.
\end{proof}
\begin{definition}[energy In Joules]
  \label{def:physics:phyx-mini-0571:phyxminiproblems-problemphyxmini0571-energyinjoules}
  \lean{PhyXMiniProblems.ProblemPhyXMini0571.energyInJoules}
  Coherent-SI readout of an energy, in joules
\end{definition}
\begin{proof}
  This is the declared model definition of energy In Joules.
\end{proof}
\begin{definition}[speed In Meters Per Second]
  \label{def:physics:phyx-mini-0571:phyxminiproblems-problemphyxmini0571-speedinmeterspersecond}
  \lean{PhyXMiniProblems.ProblemPhyXMini0571.speedInMetersPerSecond}
  \uses{def:physics:phyx-mini-0571:phyxminiproblems-problemphyxmini0571-speedquantity}
  Coherent-SI readout of a speed, in metres per second
\end{definition}
\begin{proof}
  This is the declared model definition of speed In Meters Per Second.
\end{proof}
\begin{definition}[Atom Site]
  \label{def:physics:phyx-mini-0571:phyxminiproblems-problemphyxmini0571-atomsite}
  \lean{PhyXMiniProblems.ProblemPhyXMini0571.AtomSite}
  The three labeled atomic sites in the depicted molecule
\end{definition}
\begin{proof}
  This is the declared model definition of Atom Site.
\end{proof}
\begin{definition}[Chemical Element]
  \label{def:physics:phyx-mini-0571:phyxminiproblems-problemphyxmini0571-chemicalelement}
  \lean{PhyXMiniProblems.ProblemPhyXMini0571.ChemicalElement}
  Chemical-element labels occurring in the figure
\end{definition}
\begin{proof}
  This is the declared model definition of Chemical Element.
\end{proof}
\begin{definition}[OHBond]
  \label{def:physics:phyx-mini-0571:phyxminiproblems-problemphyxmini0571-ohbond}
  \lean{PhyXMiniProblems.ProblemPhyXMini0571.OHBond}
  The two drawn oxygen-hydrogen bonds
\end{definition}
\begin{proof}
  This is the declared model definition of OHBond.
\end{proof}
\begin{definition}[hydrogen Site]
  \label{def:physics:phyx-mini-0571:phyxminiproblems-problemphyxmini0571-ohbond-hydrogensite}
  \lean{PhyXMiniProblems.ProblemPhyXMini0571.OHBond.hydrogenSite}
  \uses{def:physics:phyx-mini-0571:phyxminiproblems-problemphyxmini0571-atomsite, def:physics:phyx-mini-0571:phyxminiproblems-problemphyxmini0571-ohbond}
  The hydrogen endpoint associated with each depicted bond
\end{definition}
\begin{proof}
  This is the declared model definition of hydrogen Site.
\end{proof}
\begin{definition}[Atom Shade]
  \label{def:physics:phyx-mini-0571:phyxminiproblems-problemphyxmini0571-atomshade}
  \lean{PhyXMiniProblems.ProblemPhyXMini0571.AtomShade}
  The two shades used to distinguish the element types in the raster
\end{definition}
\begin{proof}
  This is the declared model definition of Atom Shade.
\end{proof}
\begin{definition}[Rotation Axis]
  \label{def:physics:phyx-mini-0571:phyxminiproblems-problemphyxmini0571-rotationaxis}
  \lean{PhyXMiniProblems.ProblemPhyXMini0571.RotationAxis}
  The rotation-axis role of the dashed line in the supplied diagram
\end{definition}
\begin{proof}
  This is the declared model definition of Rotation Axis.
\end{proof}
\begin{definition}[Rotational Level]
  \label{def:physics:phyx-mini-0571:phyxminiproblems-problemphyxmini0571-rotationallevel}
  \lean{PhyXMiniProblems.ProblemPhyXMini0571.RotationalLevel}
  The two rotational levels mentioned in the question
\end{definition}
\begin{proof}
  This is the declared model definition of Rotational Level.
\end{proof}
\begin{definition}[quantum Number]
  \label{def:physics:phyx-mini-0571:phyxminiproblems-problemphyxmini0571-rotationallevel-quantumnumber}
  \lean{PhyXMiniProblems.ProblemPhyXMini0571.RotationalLevel.quantumNumber}
  \uses{def:physics:phyx-mini-0571:phyxminiproblems-problemphyxmini0571-rotationallevel}
  Orbital angular-momentum quantum number attached to each named level
\end{definition}
\begin{proof}
  This is the declared model definition of quantum Number.
\end{proof}
\begin{definition}[Water Molecule Figure]
  \label{def:physics:phyx-mini-0571:phyxminiproblems-problemphyxmini0571-watermoleculefigure}
  \lean{PhyXMiniProblems.ProblemPhyXMini0571.WaterMoleculeFigure}
  \uses{def:physics:phyx-mini-0571:phyxminiproblems-problemphyxmini0571-atomsite, def:physics:phyx-mini-0571:phyxminiproblems-problemphyxmini0571-chemicalelement, def:physics:phyx-mini-0571:phyxminiproblems-problemphyxmini0571-ohbond, def:physics:phyx-mini-0571:phyxminiproblems-problemphyxmini0571-atomshade}
  Visible content of the primary molecular diagram. The numerical labels here are scalar raster readouts; the setup below separately stores the associated physical length and angle
\end{definition}
\begin{proof}
  This is the declared model definition of Water Molecule Figure.
\end{proof}
\begin{definition}[Water Molecule Rotor Setup]
  \label{def:physics:phyx-mini-0571:phyxminiproblems-problemphyxmini0571-watermoleculerotorsetup}
  \lean{PhyXMiniProblems.ProblemPhyXMini0571.WaterMoleculeRotorSetup}
  \uses{def:physics:phyx-mini-0571:phyxminiproblems-problemphyxmini0571-lengthquantity, def:physics:phyx-mini-0571:phyxminiproblems-problemphyxmini0571-massquantity, def:physics:phyx-mini-0571:phyxminiproblems-problemphyxmini0571-momentofinertiaquantity, def:physics:phyx-mini-0571:phyxminiproblems-problemphyxmini0571-actionquantity, def:physics:phyx-mini-0571:phyxminiproblems-problemphyxmini0571-speedquantity, def:physics:phyx-mini-0571:phyxminiproblems-problemphyxmini0571-atomsite, def:physics:phyx-mini-0571:phyxminiproblems-problemphyxmini0571-rotationaxis, def:physics:phyx-mini-0571:phyxminiproblems-problemphyxmini0571-rotationallevel, def:physics:phyx-mini-0571:phyxminiproblems-problemphyxmini0571-watermoleculefigure}
  Independent physical quantities of the rotor and photon. The perpendicular distances, inertia, level energies, photon energy, and photon wavelength are fields rather than definitions, and are related only by the law structures below
\end{definition}
\begin{proof}
  This is the declared model definition of Water Molecule Rotor Setup.
\end{proof}
\begin{definition}[Matches Water Rotational Excitation Scenario]
  \label{def:physics:phyx-mini-0571:phyxminiproblems-problemphyxmini0571-matcheswaterrotationalexcitationscenario}
  \lean{PhyXMiniProblems.ProblemPhyXMini0571.MatchesWaterRotationalExcitationScenario}
  \uses{def:physics:phyx-mini-0571:phyxminiproblems-problemphyxmini0571-watermoleculerotorsetup}
  The categorical scenario and the requested transition. This fixes only the roles l = 0 and l = 1, not their energy gap or the photon wavelength
\end{definition}
\begin{proof}
  This is the declared model definition of Matches Water Rotational Excitation Scenario.
\end{proof}
\begin{definition}[Matches Primary Water Molecule Figure]
  \label{def:physics:phyx-mini-0571:phyxminiproblems-problemphyxmini0571-matchesprimarywatermoleculefigure}
  \lean{PhyXMiniProblems.ProblemPhyXMini0571.MatchesPrimaryWaterMoleculeFigure}
  \uses{def:physics:phyx-mini-0571:phyxminiproblems-problemphyxmini0571-lengthinnanometers, def:physics:phyx-mini-0571:phyxminiproblems-problemphyxmini0571-watermoleculerotorsetup}
  All salient labels and qualitative geometry visible in image 571. The two physical O-H bonds share the figure's 0.0958 nm readout, and the physical angle is the radian conversion of its 105° label. No photon wavelength or answer choice occurs here
\end{definition}
\begin{proof}
  This is the declared model definition of Matches Primary Water Molecule Figure.
\end{proof}
\begin{definition}[Uses Textbook Reference Data]
  \label{def:physics:phyx-mini-0571:phyxminiproblems-problemphyxmini0571-usestextbookreferencedata}
  \lean{PhyXMiniProblems.ProblemPhyXMini0571.UsesTextbookReferenceData}
  \uses{def:physics:phyx-mini-0571:phyxminiproblems-problemphyxmini0571-massinkilograms, def:physics:phyx-mini-0571:phyxminiproblems-problemphyxmini0571-actioninjouleseconds, def:physics:phyx-mini-0571:phyxminiproblems-problemphyxmini0571-speedinmeterspersecond, def:physics:phyx-mini-0571:phyxminiproblems-problemphyxmini0571-watermoleculerotorsetup}
  Textbook reference values needed for the numerical calculation. Physlib's Constants.ℏ and DimSpeed.speedOfLight ground the two universal constants. The oxygen mass is retained even though an atom on the selected axis makes no contribution to this moment of inertia. No target wavelength appears here
\end{definition}
\begin{proof}
  This is the declared model definition of Uses Textbook Reference Data.
\end{proof}
\begin{definition}[Has Physical Water Rotor Parameters]
  \label{def:physics:phyx-mini-0571:phyxminiproblems-problemphyxmini0571-hasphysicalwaterrotorparameters}
  \lean{PhyXMiniProblems.ProblemPhyXMini0571.HasPhysicalWaterRotorParameters}
  \uses{def:physics:phyx-mini-0571:phyxminiproblems-problemphyxmini0571-lengthinmeters, def:physics:phyx-mini-0571:phyxminiproblems-problemphyxmini0571-massinkilograms, def:physics:phyx-mini-0571:phyxminiproblems-problemphyxmini0571-momentofinertiainkilogrammeterssquared, def:physics:phyx-mini-0571:phyxminiproblems-problemphyxmini0571-actioninjouleseconds, def:physics:phyx-mini-0571:phyxminiproblems-problemphyxmini0571-energyinjoules, def:physics:phyx-mini-0571:phyxminiproblems-problemphyxmini0571-speedinmeterspersecond, def:physics:phyx-mini-0571:phyxminiproblems-problemphyxmini0571-watermoleculerotorsetup}
  Positivity and physical-branch conditions for the rotor and photon
\end{definition}
\begin{proof}
  This is the declared model definition of Has Physical Water Rotor Parameters.
\end{proof}
\begin{definition}[Satisfies Depicted Bisector Geometry]
  \label{def:physics:phyx-mini-0571:phyxminiproblems-problemphyxmini0571-satisfiesdepictedbisectorgeometry}
  \lean{PhyXMiniProblems.ProblemPhyXMini0571.SatisfiesDepictedBisectorGeometry}
  \uses{def:physics:phyx-mini-0571:phyxminiproblems-problemphyxmini0571-lengthinmeters, def:physics:phyx-mini-0571:phyxminiproblems-problemphyxmini0571-watermoleculerotorsetup}
  Geometry of the dashed bisector axis. Oxygen has zero perpendicular distance to it, while each hydrogen has lever arm r sin (θ/2). This is a general geometric relation and does not contain the requested wavelength
\end{definition}
\begin{proof}
  This is the declared model definition of Satisfies Depicted Bisector Geometry.
\end{proof}
\begin{definition}[Satisfies Rigid Rotor Photon Laws]
  \label{def:physics:phyx-mini-0571:phyxminiproblems-problemphyxmini0571-satisfiesrigidrotorphotonlaws}
  \lean{PhyXMiniProblems.ProblemPhyXMini0571.SatisfiesRigidRotorPhotonLaws}
  \uses{def:physics:phyx-mini-0571:phyxminiproblems-problemphyxmini0571-lengthinmeters, def:physics:phyx-mini-0571:phyxminiproblems-problemphyxmini0571-massinkilograms, def:physics:phyx-mini-0571:phyxminiproblems-problemphyxmini0571-momentofinertiainkilogrammeterssquared, def:physics:phyx-mini-0571:phyxminiproblems-problemphyxmini0571-actioninjouleseconds, def:physics:phyx-mini-0571:phyxminiproblems-problemphyxmini0571-energyinjoules, def:physics:phyx-mini-0571:phyxminiproblems-problemphyxmini0571-speedinmeterspersecond, def:physics:phyx-mini-0571:phyxminiproblems-problemphyxmini0571-atomsite, def:physics:phyx-mini-0571:phyxminiproblems-problemphyxmini0571-rotationallevel, def:physics:phyx-mini-0571:phyxminiproblems-problemphyxmini0571-watermoleculerotorsetup}
  The governing physics used by the calculation: the point-mass moment of inertia is Σ mᵢ dᵢ² about the selected axis; a rigid rotor has E\_l = ℏ² l(l+1)/(2I); the absorbed photon energy equals the final-minus-initial level gap; and the photon obeys E λ = h c = 2πℏc. All relations are stated for independent physical fields. None specializes the wavelength to a choice value
\end{definition}
\begin{proof}
  This is the declared model definition of Satisfies Rigid Rotor Photon Laws.
\end{proof}
\begin{definition}[Answer Choice]
  \label{def:physics:phyx-mini-0571:phyxminiproblems-problemphyxmini0571-answerchoice}
  \lean{PhyXMiniProblems.ProblemPhyXMini0571.AnswerChoice}
  Labels of the four wavelength choices in the problem
\end{definition}
\begin{proof}
  This is the declared model definition of Answer Choice.
\end{proof}
\begin{definition}[wavelength Micrometers]
  \label{def:physics:phyx-mini-0571:phyxminiproblems-problemphyxmini0571-answerchoice-wavelengthmicrometers}
  \lean{PhyXMiniProblems.ProblemPhyXMini0571.AnswerChoice.wavelengthMicrometers}
  \uses{def:physics:phyx-mini-0571:phyxminiproblems-problemphyxmini0571-answerchoice}
  Wavelength printed beside each answer choice, in micrometres
\end{definition}
\begin{proof}
  This is the declared model definition of wavelength Micrometers.
\end{proof}
\begin{definition}[recorded Dataset Answer]
  \label{def:physics:phyx-mini-0571:phyxminiproblems-problemphyxmini0571-recordeddatasetanswer}
  \lean{PhyXMiniProblems.ProblemPhyXMini0571.recordedDatasetAnswer}
  \uses{def:physics:phyx-mini-0571:phyxminiproblems-problemphyxmini0571-answerchoice}
  Dataset metadata records choice C; this is not a theorem premise
\end{definition}
\begin{proof}
  This is the declared model definition of recorded Dataset Answer.
\end{proof}
\begin{definition}[answer Choice Error Micrometers]
  \label{def:physics:phyx-mini-0571:phyxminiproblems-problemphyxmini0571-answerchoiceerrormicrometers}
  \lean{PhyXMiniProblems.ProblemPhyXMini0571.answerChoiceErrorMicrometers}
  \uses{def:physics:phyx-mini-0571:phyxminiproblems-problemphyxmini0571-lengthinmicrometers, def:physics:phyx-mini-0571:phyxminiproblems-problemphyxmini0571-watermoleculerotorsetup, def:physics:phyx-mini-0571:phyxminiproblems-problemphyxmini0571-answerchoice}
  Absolute error between the physical wavelength and a displayed choice
\end{definition}
\begin{proof}
  This is the declared model definition of answer Choice Error Micrometers.
\end{proof}
\begin{definition}[Agrees With Displayed Wavelength]
  \label{def:physics:phyx-mini-0571:phyxminiproblems-problemphyxmini0571-agreeswithdisplayedwavelength}
  \lean{PhyXMiniProblems.ProblemPhyXMini0571.AgreesWithDisplayedWavelength}
  \uses{def:physics:phyx-mini-0571:phyxminiproblems-problemphyxmini0571-watermoleculerotorsetup, def:physics:phyx-mini-0571:phyxminiproblems-problemphyxmini0571-answerchoice, def:physics:phyx-mini-0571:phyxminiproblems-problemphyxmini0571-answerchoiceerrormicrometers}
  Agreement with the whole-micrometre textbook value, allowing 1 μm for the rounded atomic-mass and rigid-point-mass approximation
\end{definition}
\begin{proof}
  This is the declared model definition of Agrees With Displayed Wavelength.
\end{proof}
\begin{definition}[Is Unique Closest Wavelength Choice]
  \label{def:physics:phyx-mini-0571:phyxminiproblems-problemphyxmini0571-isuniqueclosestwavelengthchoice}
  \lean{PhyXMiniProblems.ProblemPhyXMini0571.IsUniqueClosestWavelengthChoice}
  \uses{def:physics:phyx-mini-0571:phyxminiproblems-problemphyxmini0571-watermoleculerotorsetup, def:physics:phyx-mini-0571:phyxminiproblems-problemphyxmini0571-answerchoice, def:physics:phyx-mini-0571:phyxminiproblems-problemphyxmini0571-answerchoiceerrormicrometers}
  A choice is strictly closer than every other displayed wavelength
\end{definition}
\begin{proof}
  This is the declared model definition of Is Unique Closest Wavelength Choice.
\end{proof}
% --- Archon declaration topology end ---
% --- Archon physics formalization source end ---
